\section{Problem Formulation}
\label{sec:problem}

\paragraph{Harness state.}
Let a harness
\(
  \harness=(H_{\mathrm{prompt}},H_{\mathrm{memory}},H_{\mathrm{skill}},
  H_{\mathrm{tool}},H_{\mathrm{middleware}},H_{\mathrm{route}})
\)
surround a frozen solver model.  An admissible mutation \(m\) creates one
content-addressed child \(\candidate=m(\harness)\).  It declares the changed
surface, operation, source evidence, expected behavior, and protected behavior.
Atomicity is a protocol constraint: an interaction edit is legal only when its
joint mechanism is declared and tested as such.

\paragraph{Adaptive search.}
At round \(r\), an optimizer observes an authorized view of history
\(V(\history_{<r})\), proposes candidates, and nominates at most one for the
gate.  A frozen budget \(\budget\) limits proposals, optimizer calls, solver
calls, gate queries, tokens, wall time, and failures.  The unit of replication
is a complete search run, not an individual rollout from within one adaptive
trajectory.

\paragraph{Promotion and false promotion.}
For candidate \(c\), let \(\utility(c)\) be downstream utility,
\(\cost(c)\) resource use, and \(R_k(c)\) protected-slice outcomes.  A promotion
is valid only if benefit, cost, and every protected constraint pass under a
frozen statistical rule.  A \emph{false promotion} occurs when the gate accepts
a null, harmful, inactive, non-adhered, out-of-policy, or grader-exploiting edit.
We report false promotion at the campaign level over all adaptive nominations,
not only a nominal threshold attached to one comparison.

\paragraph{Mechanism estimand.}
The proposed chain is
\[
  m \longrightarrow \activate \longrightarrow \adhere
  \longrightarrow \behavior \longrightarrow \utility,
\]
where activation means the mutated surface entered the executed request or code
path, adherence means the solver acted consistently with it, and behavior means
the preregistered intermediate outcome changed.  We do not infer an upstream
link from downstream utility.  Each link must have source-bound observations or
the result is inconclusive.

\paragraph{Threat model and claim boundary.}
The mutable plane may emit arbitrary candidate content and attempt to exploit
the task, grader, or ledger.  It may not modify split loading, grading,
promotion logic, resource accounting, sealed authorization, or historical
evidence.  The target deployment isolates candidate execution from secrets,
network, and grader state.  The current implementation establishes typed
logical, in-process boundaries only; it is not yet a hostile-code security
boundary.  Consequently, process independence and sealed-service security are
design requirements, not completed empirical or security claims.

\begin{figure}[t]
  \centering
  \resizebox{\linewidth}{!}{%
\begin{tikzpicture}[
  font=\sffamily,
  >=Latex,
  flow/.style={-Latex, line width=0.8pt},
  return/.style={-Latex, dashed, line width=0.8pt},
  box/.style={draw=black, rounded corners=2pt, align=center,
              minimum height=9mm, inner sep=4pt},
  stage/.style={box, fill=white, minimum width=21mm},
  gate/.style={box, fill=GateGreen, minimum width=20mm},
  boundary/.style={draw=black, rounded corners=4pt, inner sep=6pt,
                   line width=0.9pt},
  note/.style={align=center, font=\sffamily\footnotesize}
]
  \node[stage] (observe) {Observe\\exploration};
  \node[stage, right=4mm of observe] (diagnose) {Diagnose\\with sources};
  \node[stage, right=4mm of diagnose] (mutate) {Typed atomic\\mutation};

  \draw[flow] (observe) -- (diagnose);
  \draw[flow] (diagnose) -- (mutate);

  \begin{scope}[on background layer]
    \node[boundary, fit=(observe)(diagnose)(mutate), fill=MutableFill,
          label={[font=\bfseries\sffamily]above:Mutable evolution plane}]
          (mutable) {};
  \end{scope}

  \node[stage, below=20mm of observe, xshift=7mm] (quartet)
        {Matched quartet\\child / parent\\revert / placebo};
  \node[gate, right=4mm of quartet] (activation) {Activation};
  \node[gate, right=3mm of activation] (adherence) {Adherence};
  \node[gate, right=3mm of adherence] (behavior) {Behavior};
  \node[gate, right=3mm of behavior, minimum width=26mm] (constraints)
        {Utility + protected\\behavior + cost};

  \draw[flow] (quartet) -- (activation);
  \draw[flow] (activation) -- (adherence);
  \draw[flow] (adherence) -- (behavior);
  \draw[flow] (behavior) -- (constraints);

  \node[draw=black, dashed, rounded corners=2pt, below=7mm of quartet,
        minimum width=34mm, align=center, fill=white] (sandbox)
        {Target untrusted runtime\\no grader, targets, or secrets};
  \node[note, below=8mm of constraints, text width=58mm] (trustnote)
        {Frozen tasks, grader, schedule, evidence authorities, usage ledger,
         and campaign error budget};
  \draw[flow] (quartet.south) -- (sandbox.north);

  \begin{scope}[on background layer]
    \node[boundary,
          inner ysep=16pt,
          fit=(quartet)(activation)(adherence)(behavior)(constraints)(sandbox)(trustnote),
          fill=TrustedFill]
          (trusted) {};
  \end{scope}
  \draw[flow] (mutate.south) -- ++(0,-7mm) -| (quartet.north);
  \draw[return] (constraints.north) -- ++(0,7mm) -|
        node[pos=0.55, above, note] {redacted pass / fail / inconclusive}
        (observe.south);
  % Draw the boundary title after the dashed return so its fill masks the
  % crossing without changing the arrow's routing or meaning.
  \node[anchor=north west, font=\bfseries\sffamily, fill=TrustedFill,
        inner sep=1pt] at ($(trusted.north west)+(2mm,-1mm)$)
        {Target independent verification plane};

  \node[box, fill=SealedFill, below=10mm of trusted.south, minimum width=68mm]
        (sealed) {Target external sealed authority\\all-cell barrier $\rightarrow$
        commit registry + champions $\rightarrow$ one atomic release};
  \draw[flow] (trusted.south) --
        node[right, note] {all search cells close} (sealed.north);
  \node[note, right=3mm of sealed.east, align=left]
        {Never visible to\\optimizer or candidate};
\end{tikzpicture}%
}

  \caption{Mechanism and authority boundary.  The mutable optimizer can edit
  only declared harness surfaces.  Promotion belongs to a separate authority
  that executes matched interventions and fails closed along the mechanism
  chain.  Sealed tasks are released once, only after all search cells close.
  Solid arrows carry typed artifacts; the dashed arrow carries a redacted
  decision, never sealed item feedback.  Shape and labels preserve the meaning
  in grayscale.}
  \label{fig:mechanism-boundary}
\end{figure}
